\setcounter{exercisegroup}{0}
\refstepcounter{exercisegroup}
\section*{第 \S\arabic{exercisegroup} 节习题}
\phantomsection
\addcontentsline{toc}{section}{第 \protect\S\arabic{exercisegroup} 节习题}

\begin{enumerate}
\def\labelenumi{\arabic{enumi})}
\item
  设 \(\mathbf{f}\) 是 \(\mathbf{R}\) 到 \(\mathbf{R}\) 上的一个具有单位元 \(\mathbf{e}\) 的完备赋范代数 E 中的连续且可导的映射；假设 \(\mathbf{f}(0) = \mathbf{e}\) ，并且恒有 \(\mathbf{f}'(x) = \mathbf{f}(x)\mathbf{c}\) ，其中 \(\mathbf{c}\) 是 E 的可逆元。试证：\(\mathbf{f}\) 是加法群 \(\mathbf{R}\) 到 E 的可逆元所成的乘法群中的同态。（考虑包含 0 的最大开区间 I ，使得 \(\mathbf{f}(z)\) 对一切 \(z \in I\) 都可逆，并证明当 \(x, y\) 与 \(x + y\) 属于 I 时，有 \(\mathbf{f}(x + y)(\mathbf{f}(x)\mathbf{f}(y))^{-1} = \mathbf{e}\) ，其中保持 \(x\) 固定而让 \(y\) 变动；然后推出 \(I = R\) 。）
\item
  \begin{enumerate}
  \def\labelenumii{\alph{enumii})}
  \tightlist
  \item
    为使函数 \((1 + 1/x)^{x+p}\) 对 x \textgreater{} 0 递减（相应地，递增），必须且只需 \(p \geqslant \frac{1}{2}\) （相应地，\(p \leqslant 0\) ）；当 \(0 < p < \frac{1}{2}\) 时，该函数在区间 {]}0, \(x_{0}[\) 上递减，而在 {]} \(x_{0}, +\infty[\) 上递增。在每种情形下，当 x 趋于 \(+\infty\) 时该函数都趋于 e 。
  \end{enumerate}
\end{enumerate}

\begin{enumerate}
\def\labelenumi{\alph{enumi})}
\setcounter{enumi}{1}
\tightlist
\item
  以同样的方式研究函数 \((1 - 1 / x)^{x - p}\) 、\((1 + 1 / x)^{x}(1 + p / x)\) 与 \((1 + p / x)^{x + 1}\) 对 \(x > 0\) 的性态。
\end{enumerate}

\begin{enumerate}
\def\labelenumi{\arabic{enumi})}
\setcounter{enumi}{2}
\tightlist
\item
  \begin{enumerate}
  \def\labelenumii{\alph{enumii})}
  \tightlist
  \item
    试证：\((1+x)^{\alpha} \geqslant 1 + \alpha x\) 对 \(\alpha \geqslant 1\) 与 \(x \geqslant 0\) 成立，而 \((1+x)^{\alpha} \leqslant 1 + \alpha x\) 对 \(0 \leqslant \alpha \leqslant 1\) 与 \(x \geqslant 0\) 成立。
  \end{enumerate}
\end{enumerate}

\begin{enumerate}
\def\labelenumi{\alph{enumi})}
\setcounter{enumi}{1}
\item
  试证：\((1 - x)^{\alpha} \leqslant \frac{1}{1 + \alpha x}\) 对 \(0 \leqslant x \leqslant 1\) 与 \(\alpha \geqslant 0\) 成立。
\item
  当 \(0 \leqslant \alpha \leqslant 1\) 时，试证：对 \(x \geqslant 0\) 与 \(y \geqslant 0\) ，有 \(x^{\alpha} y^{1 - \alpha} \leqslant \alpha x + by\) ，对每一对 \(>0\) 且满足 \(a^{\alpha} b^{1 - \alpha} = \alpha (1 - \alpha)^{1 - \alpha}\) 的数成立。
\item
  由 c) 推出：若 f 、g 是区间 I 上两个 \(\geqslant 0\) 的规则函数，使得积分 \(\int_{I} f(t) dt\) 与 \(\int_{I} g(t) dt\) 收敛，则积分 \(\int_{I} (f(t))^{\alpha}(g(t))^{1-\alpha} dt\) 收敛，并且
\end{enumerate}

\[
\int_ {I} (f (t)) ^ {\alpha} (g (t)) ^ {1 - \alpha} d t \leqslant a \int_ {I} f (t) d t + b \int_ {I} g (t) d t.
\]

由此推出赫尔德不等式

\[
\int_ {I} (f (t)) ^ {\alpha} (g (t)) ^ {1 - \alpha} d t \leqslant \left(\int_ {I} f (t) d t\right) ^ {\alpha} \left(\int_ {I} g (t) d t\right) ^ {1 - \alpha}
\]

（当 \(\alpha = \frac{1}{2}\) 时，它称为“柯西–布尼亚科夫斯基–施瓦茨不等式”）。

\begin{enumerate}
\def\labelenumi{\arabic{enumi})}
\setcounter{enumi}{3}
\tightlist
\item
  由柯西–施瓦茨不等式推出：对每个函数 \(u\) （它是 \([a, b]\) 上某个规则函数的原函数），以及每个连续函数 \(h > 0\) （在 \([a, b]\) 上），有
\end{enumerate}

\[
\left| u (b) ^ {2} - u (a) ^ {2} \right| \leqslant 2 \left(\int_ {a} ^ {b} u ^ {\prime 2} (t) h (t) d t\right) ^ {1 / 2} \left(\int_ {a} ^ {b} \frac {u ^ {2} (t)}{h (t)} d t\right) ^ {1 / 2}
\]

只要右端的两个积分收敛。

取 \(a = 0\) ，\(b = \pi / 2\) ，\(h(t) = 1\) 以及

\[
u (t) = c _ {1} \cos t + c _ {2} \cos 3 t + \dots + c _ {n} \cos (2 n - 1) t,
\]

其中对每个 j 有 \(c_{j} \geqslant 0\) ，由此推出卡尔森不等式

\[
\sum_ {j} c _ {j} \leqslant \sqrt {\pi} \left(\sum_ {j} c _ {j} ^ {2}\right) ^ {1 / 4} \left(\sum_ {j} \left(j - \frac {1}{2}\right) ^ {2} c _ {j} ^ {2}\right) ^ {1 / 4}.
\]

\begin{enumerate}
\def\labelenumi{\arabic{enumi})}
\setcounter{enumi}{4}
\tightlist
\item
  对 \(x > 0\) 及任意实数 \(y\) ，试证
\end{enumerate}

\[
x y \leqslant \log x + e ^ {y - 1}
\]

（cf.~II, p.~82, exerc. 15）；在什么情形下两边相等？

\begin{enumerate}
\def\labelenumi{\arabic{enumi})}
\setcounter{enumi}{5}
\tightlist
\item
  用 \(A_{n}\) 与 \(G_{n}\) 表示由正数组成的序列 \((a_{k})\) （ \(k \geqslant 1\) ）的前 n 个数的通常算术平均与几何平均。试证
\end{enumerate}

\[
n \left(\mathrm{A} _ {n} - \mathrm{G} _ {n}\right) \leqslant (n + 1) \left(\mathrm{A} _ {n + 1} - \mathrm{G} _ {n + 1}\right)
\]

等号仅当 \(a_{n+1} = G_{n}\) 时成立（令 \(a_{n+1} = x^{n+1}\) ，\(G_{n} = y^{n+1}\) ）。

\begin{enumerate}
\def\labelenumi{\arabic{enumi})}
\setcounter{enumi}{6}
\tightlist
\item
  用习题 6 的记号，试证：若令 \(x_{i} = a_{i} / \mathrm{A}_{n}\) ，则关系式 \(\mathrm{G}_n \geqslant (1 - \alpha)\mathrm{A}_n\) （对 \(0 \leqslant \alpha < 1\) ）蕴含：对 \(1 \leqslant i \leqslant n\) ，
\end{enumerate}

\[
x _ {i} \left(1 - \frac {x _ {i} - 1}{n - 1}\right) ^ {n - 1} \geqslant (1 - \alpha) ^ {n}.
\]

推出：对每个指标 \(i\) ，有 \(1 + x' \leqslant x_i \leqslant 1 + x''\) ，其中 \(x'\) 是方程的 \(\leqslant 0\) 的根，\(x''\) 是方程的 \(\geqslant 0\) 的根，这里方程指 \((1 + x)e^{-x} = (1 - \alpha)^n\) （cf.~III, p.~115, exerc. 2）。

\begin{enumerate}
\def\labelenumi{\arabic{enumi})}
\setcounter{enumi}{7}
\tightlist
\item
  考虑一列集合 \(D_{k}\) （ \(k \geqslant 1\) ），并对每个 k 考虑两个实值函数 \((x_{1}, \ldots, x_{k}) \mapsto f_{k}(x_{1}, \ldots, x_{k}), \quad (x_{1}, \ldots, x_{k}) \mapsto g_{k}(x_{1}, \ldots, x_{k})\) ，其中 \((x_{1}, \ldots, x_{k}) \in D_{1} \times D_{2} \times \cdots \times D_{k}\) 。假设对每个 k 都存在一个定义在 R 上的实函数 \(F_{k}\) ，使得对每个 \(\mu \in R\) 及任意的 \(a_{1} \in D_{1}, \ldots, a_{k-1} \in D_{k-1}\) ，有
\end{enumerate}

\[
\sup _ {x _ {k} \in \mathrm{D} _ {k}} \left(\mu f _ {k} (a _ {1}, \dots , a _ {k - 1}, x _ {k}) - g _ {k} (a _ {1}, \dots , a _ {k - 1}, x _ {k})\right) = \mathrm{F} _ {k} (\mu) f _ {k - 1} (a _ {1}, \dots , a _ {k - 1}).
\]

（按约定取 \(f_0 = 1\) 。）试证：若实数序列 \((\mu_k)\) 与 \((\delta_k)\) 满足条件 \(\mathrm{F}_1(\mu_1) = 0\) 与 \(\mathrm{F}_k(\mu_k) = \delta_k\) ，则有不等式

\[
\left(\mu_ {1} - \delta_ {2}\right) f _ {1} + \left(\mu_ {2} - \delta_ {3}\right) f _ {2} + \dots + \left(\mu_ {n - 1} - \delta_ {n}\right) f _ {n - 1} + \mu_ {n} f _ {n} \leqslant g _ {1} + g _ {2} + \dots + g _ {n}
\]

它对一切 n 及一切 \(x_{k} \in D_{k}\) 成立。

\begin{enumerate}
\def\labelenumi{\arabic{enumi})}
\setcounter{enumi}{8}
\tightlist
\item
  用习题 6 的记号，试证
\end{enumerate}

\[
\mu_ {1} \mathrm{G} _ {1} + \mu_ {2} \mathrm{G} _ {2} + \dots + \mu_ {n} \mathrm{G} _ {n} \leqslant \lambda_ {1} a _ {1} + \lambda_ {2} a _ {2} + \dots + \lambda_ {n} a _ {n}
\]

只要 \(\lambda_{k}\geqslant 0\) 对每个 \(k\) 成立，并且

\[
\mu_ {k} = k \left(\left(\lambda_ {k} \beta_ {k}\right) ^ {1 / k} - \beta_ {k + 1} ^ {1 / k}\right) \quad \text { for } 1 \leqslant k \leqslant n
\]

其中 \(\beta_{k}\geqslant 0\) 对一切 \(k\) 成立，\(\beta_{1}\leqslant 1\) 且 \(\beta_{n + 1} = 0\) （习题 8 的方法）。

特别地，适当选取 \(\lambda_{k}\) 与 \(\beta_{k}\) ，便有

\[
\mathrm{G} _ {1} + \mathrm{G} _ {2} + \dots + \mathrm{G} _ {n} + n \mathrm{G} _ {n} \leqslant 2 a _ {1} + \left(\frac {3}{2}\right) ^ {2} a _ {2} + \dots + \left(\frac {n + 1}{n}\right) ^ {n} a _ {n}
\]

而特别地有（卡莱曼不等式）

\[
\begin{array}{c} \mathrm{G} _ {1} + \mathrm{G} _ {2} + \dots + \mathrm{G} _ {n} <   e (a _ {1} + a _ {2} + \dots + a _ {n}); \\ e \mathrm{A} _ {n} \geqslant \Gamma_ {n} e ^ {\mathrm{G} _ {n} / \Gamma_ {n}} \end{array}
\]

其中 \(\Gamma_{n}=(\mathrm{G}_{1}+\mathrm{G}_{2}+\cdots+\mathrm{G}_{n})/n\) （卡莱曼不等式的加强）；

\[
\frac {\mathrm{A} _ {n}}{\mathrm{G} _ {n}} \geqslant \frac {1}{2} \left(\frac {\Gamma_ {n}}{\mathrm{G} _ {n}} + \frac {\mathrm{G} _ {n}}{\Gamma_ {n}}\right)
\]

若序列 \((a_{k})\) 递增；

\[
n \mathrm{G} _ {n} - m \mathrm{G} _ {m} \leqslant a _ {m + 1} + a _ {m + 2} + \dots + a _ {n}.
\]

\begin{enumerate}
\def\labelenumi{\arabic{enumi})}
\setcounter{enumi}{9}
\tightlist
\item
  设序列 \((a_{k})_{k\geqslant1}\) 由 \textgreater0 的数组成，我们令 \(s_{n}=nA_{n}=a_{1}+\cdots+a_{n}\) 。对 p\textless1 与 \(\lambda_{k}>0\) ，有
\end{enumerate}

\[
\mu_ {1} s _ {1} ^ {1 / p} + \mu_ {2} s _ {2} ^ {1 / p} + \dots + \mu_ {n} s _ {n} ^ {1 / p} \leqslant \lambda_ {1} a _ {1} ^ {1 / p} + \lambda_ {2} a _ {2} ^ {1 / p} + \dots + \lambda_ {n} a _ {n} ^ {1 / p}
\]

只要 \(\mu_1 \geqslant \lambda_1\) ，且对 \(k \geqslant 2\) 有 \(\mu_k = \left( \lambda_k^q + \delta_k^q \right)^{1/\delta} - \delta_{k+1}\) ，其中 \(q = p/(1-p)\) ，\(\delta_k \geqslant 0\) 对每个 \(k\) 成立，且 \(\delta_{n+1} = 0\) （习题 8 的方法）。特别地，有不等式

\[
\sum_ {k = 1} ^ {n} \left((k - p) ^ {1 / q} - k ^ {1 / q}\right) s _ {k} ^ {1 / q} + n ^ {1 / q} s _ {n} ^ {1 / q} \leqslant (1 - p) ^ {1 / q} \sum_ {k = 1} ^ {n} a _ {k} ^ {1 / q}
\]

它们蕴含（令 \(\mathrm{A}_n = s_n / n\) ）

\[
\sum_ {k = 1} ^ {n} \mathrm{A} _ {k} ^ {1 / p} + \frac {n}{1 - p} \mathrm{A} _ {n} ^ {1 / p} \leqslant (1 - p) ^ {- 1 / p} \sum_ {k = 1} ^ {n} a _ {k} ^ {1 / p}
\]

（加强的哈代不等式）；

\[
\frac {1}{q} \sum_ {k = 1} ^ {n} \left(b _ {k} - 1\right) A _ {k} ^ {1 / p} + n A _ {n} ^ {1 / p} \leqslant \sum_ {k = 1} ^ {n} a _ {k} ^ {1 / p} b _ {k} ^ {1 / p}
\]

只要 \(b_{k} \geqslant 1\) 对一切 \(k\) 成立；

\[
\alpha \left(a _ {1} ^ {2} + a _ {2} ^ {2} + \dots + a _ {n} ^ {2}\right) \leqslant a _ {1} ^ {2} + \left(a _ {2} - a _ {1}\right) ^ {2} + \dots + \left(a _ {n} - a _ {n - 1}\right) ^ {2} + \beta a _ {n} ^ {2}
\]

其中 \(\alpha = 2(1 - \cos \theta)\) 且 \(\beta = 1 - \frac{\sin(n + 1)\theta}{\sin n\theta}\) ，\(\theta\) 满足 \(0\leqslant \theta < \frac{\pi}{n}\) 。

\(\theta = \frac{\pi}{n + 1}\) 的情形；\(\theta = \frac{\pi}{2n + 1}\) 的情形。

\begin{enumerate}
\def\labelenumi{\arabic{enumi})}
\setcounter{enumi}{10}
\tightlist
\item
  设 \(a_{ij}\) （ \(1 \leqslant i \leqslant n\) ， \(1 \leqslant j \leqslant n\) ）是一些 \(\geqslant 0\) 的数，使得对每个指标 \(i\) 有 \(\sum_{j=1}^{n} a_{ij} = 1\) ，且对每个指标 \(j\) 有 \(\sum_{i=1}^{n} a_{ij} = 1\) 。设 \(x_i\) （ \(1 \leqslant i \leqslant n\) ）是 \(n\) 个数
\end{enumerate}

\[
> 0;
\]

\[
y _ {i} = a _ {i 1} x _ {1} + a _ {i 2} x _ {2} + \dots + a _ {i n} x _ {n} \quad (1 \leqslant i \leqslant n).
\]

试证 \(y_{1}y_{2}\ldots y_{n}\geqslant x_{1}x_{2}\ldots x_{n}\) （对每个指标 i 估计 \(\log y_{i}\) 的下界）。

\begin{enumerate}
\def\labelenumi{\arabic{enumi})}
\setcounter{enumi}{11}
\tightlist
\item
  \begin{enumerate}
  \def\labelenumii{\alph{enumii})}
  \tightlist
  \item
    设 \(\sum_{i,j} c_{ij} x_i \bar{x}_j\) （其中 \(c_{ji} = \bar{c}_{ij}\) ）是一个非退化的正埃尔米特型，并且
  \end{enumerate}
\end{enumerate}

\(\Delta\) 是它的行列式；试证 \(\Delta \leqslant c_{11}c_{22}\ldots c_{nn}\) （用该埃尔米特型的特征值表示 \(\Delta\) 与诸 \(c_{ii}\) ，并利用命题 2）。

\begin{enumerate}
\def\labelenumi{\alph{enumi})}
\setcounter{enumi}{1}
\tightlist
\item
  推出：若 \((a_{ij})\) 是以任意复数作为元素的 n 阶方阵，而 \(\Delta\) 是它的行列式，则有（“阿达马不等式”）
\end{enumerate}

\[
\left| \Delta \right| ^ {2} \leqslant \left(\sum_ {j = 1} ^ {n} \left| a _ {1 j} \right| ^ {2}\right) \left(\sum_ {j = 1} ^ {n} \left| a _ {2 j} \right| ^ {2}\right) \dots \left(\sum_ {j = 1} ^ {n} \left| a _ {n j} \right| ^ {2}\right)
\]

等号仅当右端诸项中有一项为零，或当对所有互异的指标 \(h, k\)

\[
a _ {h 1} \bar {a} _ {k 1} + a _ {h 2} \bar {a} _ {k 2} + \dots + a _ {h n} \bar {a} _ {k n} = 0
\]

（把矩阵 \((a_{ij})\) 乘以它的转置矩阵的共轭）。

\begin{enumerate}
\def\labelenumi{\arabic{enumi})}
\setcounter{enumi}{12}
\tightlist
\item
  若 \(x, y, a, b\) 均 \(> 0\) ，试证
\end{enumerate}

\[
x \log \frac {x}{a} + y \log \frac {y}{b} \geqslant (x + y) \log \frac {x + y}{a + b}
\]

等号仅当 x/a = y/b 时成立。

\begin{enumerate}
\def\labelenumi{\arabic{enumi})}
\setcounter{enumi}{13}
\tightlist
\item
  设 \(a\) 是满足 \(0 < a < \pi / 2\) 的实数。试证函数
\end{enumerate}

\[
\frac {\frac {\tan x}{x} - \frac {\tan a}{a}}{x \tan x - a \tan a}
\]

在区间 \(a, \pi/2\) 上严格递增（cf.~I, p.~39, exerc. 11）。

\begin{enumerate}
\def\labelenumi{\arabic{enumi})}
\setcounter{enumi}{14}
\item
  设 \(u\) 与 \(v\) 是 \(x\) 的两个多项式（实系数），使得恒有 \(\sqrt{1 - u^2} = v\sqrt{1 - x^2}\) ；试证：若 \(n\) 是 \(u\) 的次数，则 \(u' = nv\) ；并推出 \(u(x) = \cos(n\operatorname{Arc}\cos x)\) 。
\item
  试证（对 \(n\) 用归纳法）
\end{enumerate}

\[
\mathrm{D} ^ {n} \left(\operatorname{Arc} \tan x\right) = (- 1) ^ {n - 1} \frac {(n - 1) !}{(1 + x ^ {2}) ^ {n / 2}} \sin \left(n \operatorname{Arc} \tan \frac {1}{x}\right).
\]

¶ 17) 设 \(f\) 是定义在开区间 \(\mathrm{I}\subset \mathbf{R}\) 上的实函数，使得对任意一组三点 \(x_{1}, x_{2}, x_{3}\) （属于 I 且满足 \(x_{1} < x_{2} < x_{3} < x_{1} + \pi\) ），有

\[
f \left(x _ {1}\right) \sin \left(x _ {3} - x _ {2}\right) + f \left(x _ {2}\right) \sin \left(x _ {1} - x _ {3}\right) + f \left(x _ {3}\right) \sin \left(x _ {2} - x _ {1}\right) \geqslant 0.\tag{1}
\]

试证：

\begin{enumerate}
\def\labelenumi{\alph{enumi})}
\tightlist
\item
  \(f\) 在 I 的每一点处连续，并且在 I 的每一点处有有限的右导数与有限的左导数；此外，
\end{enumerate}

\[
f (x) \cos (x - y) - f _ {d} ^ {\prime} (x) \sin (x - y) \leqslant f (y)\tag{2}
\]

对 I 中每一对点 \(x, y\) （满足 \(|x - y| \leqslant \pi\) ）成立；把 \(f_d'\) 换成 \(f_g'\) ，(2) 也有一个类似的关系式；最后，有 \(f_g'(x) \leqslant f_d'(x)\) 对每个 \(x \in I\) 成立。（为证 (2) ，在 (1) 中让 \(x_2\) 趋于 \(x_1\) 而保持 \(x_3\) 固定，从而得到 \(\limsup_{x_2 \to x_1} \frac{f(x_2) - f(x_1)}{x_2 - x_1}\) 的一个上界；然后，在所得的不等式中让 \(x_3\) 趋于 \(x_1\) ；由此推出 \(f_d'(x)\) 的存在性以及 \(y > x\) 时的不等式 (2) ；对其余部分同样处理。）

\begin{enumerate}
\def\labelenumi{\alph{enumi})}
\setcounter{enumi}{1}
\item
  反之，若 (2) 对 I 中每一对点 \(x, y\) （满足 \(|x - y| \leqslant \pi\) ）成立，则 \(f\) 在 I 上满足 (1) （考虑差 \(\frac{f(x_2)}{\sin(x_3 - x_2)} - \frac{f(x_1)}{\sin(x_3 - x_1)}\) ）。
\item
  若 \(f\) 在 \(I\) 上有二阶导数，则 (1) 等价于条件
\end{enumerate}

\[
f (x) + f ^ {\prime \prime} (x) \geqslant 0\tag{3}
\]

它对每个 \(x\in \mathrm{I}\) 成立。

*试解释这些结果：考虑由 \(x = \frac{1}{f(t)} \cos t\) ，\(y = \frac{1}{f(t)} \sin t\) 定义的平面曲线（“关于原点的凸性”）。*

\begin{enumerate}
\def\labelenumi{\arabic{enumi})}
\setcounter{enumi}{17}
\tightlist
\item
  \begin{enumerate}
  \def\labelenumii{\alph{enumii})}
  \tightlist
  \item
    试证：在从 \(\mathbf{R}\) 到 \(\mathbf{C}\) 的映射所成的向量空间中，形如 \(x^n e^{\alpha x}\) （ \(n\) 为整数，\(\alpha\) 为任意复数）的不同函数构成一个线性无关组（用反证法论证：考虑这些函数之间系数 \(\neq 0\) 的个数尽可能少的一个关系式，然后对它求导）。
  \end{enumerate}
\end{enumerate}

\begin{enumerate}
\def\labelenumi{\alph{enumi})}
\setcounter{enumi}{1}
\tightlist
\item
  设 \(f(\mathrm{X}_1, \mathrm{X}_2, \ldots, \mathrm{X}_m)\) 是 \(m\) 个未定元的多项式（复系数），使得当用函数 \(\cos\left(\sum_{k=1}^{n} p_{jk} x_k\right)\) 代换 \(\mathrm{X}_j\) （对 \(1 \leqslant j \leqslant r\) ），并用函数 \(\sin\left(\sum_{k=1}^{n} p_{jk} x_k\right)\) （对 \(r + 1 \leqslant j \leqslant m\) ；其中 \(p_{jk}\) 为实数，\(x_k\) 为实数）时，得到关于 \(x_k\) 恒为零的函数；试证：当给 \(x_k\) 以任意的复数值时，同一个恒等式仍然成立（利用 \(a\) ）。
\end{enumerate}

¶ 19) 我们知道（A, IX, §10, exerc. 2），在复平面 \(\mathbf{C}^2\) 中，定向直线的角所成的群 A 同构于正交群 \(\mathbf{O}_2(\mathbf{C})\) ，其中的典范同构把角 \(\theta\) 对应于转角为 \(\theta\) 的旋转；通过这个同构，把 \(\mathbf{O}_2(\mathbf{C})\) （被视为矩阵空间 \(\mathbf{M}_2(\mathbf{C})\) （ \(\mathbf{C}\) 上 2 阶矩阵所成的空间）的子空间）的拓扑转移到群 A 上，从而使 A 成为一个局部紧拓扑群；

此外，映射 \(\theta \mapsto \cos \theta + i\sin \theta\) 是拓扑群 A 到乘法群 \(\mathbf{C}^*\) （由 \(\neq 0\) 的复数组成）上的同构，其逆同构由公式 \(\cos \theta = \frac{1}{2}(z + 1 / z)\) ，\(\sin \theta = \frac{1}{2i} (z - 1 / z)\) 定义。由这些关系式推出：加法群 C 到 A 上的每个连续同态 \(z \mapsto \varphi(z)\) ，只要复函数 \(\cos \varphi(z)\) ，\(\sin \varphi(z)\) ) 在 C 上可导，就由关系式 \(\cos \varphi(z) = \cos az\) ，\(\sin \varphi(z) = \sin az\) 定义（ \(a\) 为复数）。

\begin{enumerate}
\def\labelenumi{\arabic{enumi})}
\setcounter{enumi}{19}
\item
  设 D 是 C 的下述子集：由 \(-\pi < \mathcal{R}(z) \leqslant \pi\) ，\(\mathcal{I}(z) > 0\) 所定义的集与线段 \(\mathcal{I}(z) = 0\) ，\(0 \leqslant \mathcal{R}(z) \leqslant \pi\) 的并。试证：函数 \(\cos z\) 在 D 上的限制是 D 到 C 上的一个双射；\(\cos z\) 在 D 的内部的限制是这个开集到 C 中半直线 \(y = 0\) ，\(x \leqslant 1\) 的余集上的一个同胚。
\item
  设 \(f\) 与 \(g\) 是两个互素的多项式（复系数），\(f\) 的次数严格小于 \(g\) 的次数。设 \(p\) 是 \(g\) 与它的导数 \(g'\) 的最大公约数，\(q\) 是 \(g\) 除以 \(p\) 的商；试证：存在两个唯一确定的多项式 \(u, v\) ，它们的次数分别严格小于 \(p\) 与 \(q\) 的次数，使得
\end{enumerate}

\[
{\frac {f}{g}} = \mathrm{D} \left({\frac {u}{p}}\right) + {\frac {v}{q}}.
\]

推出：\(u\) 与 \(v\) 的系数属于（ \(\mathbf{Q}\) 上的）包含 \(f\) 与 \(g\) 的系数且含于 \(\mathbf{C}\) 中的最小域。

\begin{enumerate}
\def\labelenumi{\arabic{enumi})}
\setcounter{enumi}{21}
\tightlist
\item
  设 \(f\) 与 \(g\) 是两个互素的多项式（复系数），\(f\) 的次数严格小于 \(g\) 的次数。设 \(K\) 是 \(C\) 的一个子域，它包含 \(f\) 与 \(g\) 的系数，且使得 \(g\) 在 \(K\) 上不可约。为使 \(f / g\) 有一个形如 \(\sum_{i} a_i \log u_i\) 的原函数，其中 \(a_i\) 是属于 \(K\) 的常数，而 \(u_i\)
\end{enumerate}

是 K 上的不可约多项式，必须且只需 \(f = cg'\) ，其中 c 是 K 中的常数。

\begin{enumerate}
\def\labelenumi{\arabic{enumi})}
\setcounter{enumi}{22}
\item
  若 \(f(x, y)\) 是 x 、y 的任意一个多项式（复系数），试证：求 \(f(x, \log x)\) 与 \(f(x, \operatorname{Arc} \sin x)\) 的原函数都可以归结为求某个有理函数的原函数。
\item
  试证：可以把求 \((ax + b)^{p} x^{q}\) （p 与 q 为有理数）的原函数归结为求某个有理函数的原函数，只要数 p 、q 、\(p + q\) 之一是整数（正的或负的）。
\end{enumerate}

* 25) 开圆盘 \(\Delta\) （ \(\mathbf{C}\) 中的）上的亚纯函数构成一个域 \(\mathrm{M}(\Delta)\) 。\(\mathrm{M}(\Delta)\) 的一个子域 F ，若满足 \(u \in \mathrm{F}\) 蕴含 \(\mathrm{D}u \in \mathrm{F}\) ，就称为 \(\mathrm{M}(\Delta)\) 的一个微分子域。

\begin{enumerate}
\def\labelenumi{\alph{enumi})}
\tightlist
\item
  对每个多项式 \(\mathrm{P}(\mathrm{X})=\mathrm{X}^{m}+a_{1}\mathrm{X}^{m-1}+\cdots+a_{m}\) （系数属于 \(\mathrm{M}(\Delta)\) ），试证：存在一个开圆盘 \(\Delta_{1}\subset\Delta\) （一般说来它与 \(\Delta\) 不同心）和一个函数 \(f\in\mathrm{M}(\Delta_{1})\) ，满足
\end{enumerate}

\[
\Big (f (z) \Big) ^ {m} + a _ {1} (z) \Big (f (z) \Big) ^ {m - 1} + \dots + a _ {m} (z) = 0
\]

在点 \(z \in \Delta_1\) （使 \(f\) 与诸 \(a_j\) 全纯者）处都成立。若 \(F\) 是 \(\mathbf{M}(\Delta)\) 的包含诸 \(a_j\) 的子域，则称 \(\mathbf{M}(\Delta_1)\) 中由 \(f\) 与 F 中的函数在 \(\Delta_{1}\) 上的限制生成的子域为把 P 的根 f 添加到 F 上而得到的域 \(\mathrm{F}(f)\) 。这种措辞上的滥用不会引起混淆，因为限制映射 \(g \mapsto g|\Delta_{1}\) （它把 F 映到子域 \(\mathrm{M}(\Delta_{1})\) 上）是单射。若 F 是微分域，则 \(\mathrm{F}(f)\) 也是微分域。

\begin{enumerate}
\def\labelenumi{\alph{enumi})}
\setcounter{enumi}{1}
\item
  对每个函数 \(a \in \mathbf{M}(\Delta)\) ，试证：存在一个开圆盘 \(\Delta_{2} \subset \Delta\) ，使得有一个函数 \(g \in \mathbf{M}(\Delta_{2})\) 满足关系 \(g'(z) = a(z)\) （在使 g 与 a 全纯的每一点 \(z \in \Delta_{2}\) 处）。若 F 是 \(\mathbf{M}(\Delta)\) 的包含 a 的子域，则称子域 \(\mathrm{F}(g)\) （ \(\mathbf{M}(\Delta_{2})\) 中由 g 与 F 中的函数在 \(\Delta_{2}\) 上的限制生成者）为把原函数 \(\int a dz\) 添加到 F 上而得到的域。若 F 是微分域，则 \(\mathrm{F}(g)\) 也是微分域。
\item
  对每个函数 \(b \in \mathrm{M}(\Delta)\) ，试证：存在一个开圆盘 \(\Delta_{3} \subset \Delta\) ，使得有一个 \(h \in \mathrm{M}(\Delta_{3})\) 满足关系 \(h'(z) = b(z)h(z)\) （在每一点 \(z \in \Delta_{3}\) 处，只要 h 与 b 全纯且 \(\neq 0\) ）。若 F 是 \(\mathrm{M}(\Delta)\) 的包含 b 的子域，则称子域 \(\mathrm{F}(h)\) （ \(\mathrm{M}(\Delta_{3})\) 中由 h 与 F 中的函数在 \(\Delta_{3}\) 上的限制生成者）为把原函数的指数式 \(\exp\left(\int bdz\right)\) 添加到 F 上而得到的域。若 F 是微分域，则 \(\mathrm{F}(h)\) 也是微分域。
\item
  对 M( \(\Delta\) ) 的每个子域 F 及每个系数属于 F 的多项式 P(X) = X \(^{m}\) + a \(_{1}\) X \(^{m-1}\) + \(\cdots\) + a \(_{m}\) ，存在一个域 K ，它是把一些多项式的根逐次添加到 F 上而得到的，使得 K 是 F 的一个伽罗瓦扩张，并且有 P(X) = (X - c \(_{1}\) ) \ldots{} (X - c \(_{m}\) ) ，其中 c \(_{i}\) 是属于 K 的亚纯函数（在适当的圆盘上）。若 F 是微分域，则有 ( \(\sigma.g\) )' = \(\sigma.g'\) ，对每个 g ∈ K 及 K 在 F 上的伽罗瓦群的每个元素 \(\sigma\) 成立。*
\end{enumerate}

* 26) a) 设 \(\mathrm{F} \subset \mathrm{M}(\Delta)\) 是一个微分域，\(\mathbf{K}\) 是 \(\mathbf{F}\) 的一个有限伽罗瓦扩张，且是某个 \(\mathrm{M}(\Delta_1)\) 的子域。设 \(t\) 是 \(\mathbf{F}\) 中某个函数的原函数或原函数的指数式；试证：若 \(t\) 在 \(\mathbf{F}\) 上超越，则不存在函数 \(u \in \mathbf{K}\) 使得 \(t' = u'\) 。（分别考虑 \(t' = a \in \mathbf{F}\) 与 \(t' = bt\) （ \(b \in \mathbf{F}\) ）两种情形；考虑变换 \(\sigma.u\) （ \(u\) 在 \(\mathbf{K}\) 于 \(\mathbf{F}\) 上的伽罗瓦群作用下）及其导数 \(\sigma.u'\) ，由此得出矛盾；在第一种情形，证明将有 \(t' = c'\) （对某个 \(c \in \mathbf{F}\) ）；在第二种情形，令 \(\mathrm{N} = [\mathrm{K}: \mathrm{F}]\) ，将有 \((t^{N}/c)' = 0\) （对某个 \(c \in \mathrm{F}\) ）。）

\begin{enumerate}
\def\labelenumi{\alph{enumi})}
\setcounter{enumi}{1}
\tightlist
\item
  设 t 在 F 上超越且 \(t' = bt\) （ \(b \in F\) ）；试证：K 中不存在任何 \(c \neq 0\) 使得 \((ct^{m})' = 0\) （同样的方法）。*
\end{enumerate}

* 27) 设 \(\mathrm{F} \subset \mathrm{M}(\Delta)\) 是一个包含 \(\mathbf{C}\) 的微分域，\(t\) 是 \(\mathrm{F}\) 中某个函数的原函数或原函数的指数式，并设 \(t\) 在 \(\mathrm{F}\) 上超越。又设 \(c_1, \ldots, c_n\) 是 \(\mathrm{F}\) 中在有理数域 \(\mathbf{Q}\) 上线性无关的元素；那么，若 \(u_1, \ldots, u_n\) 与 \(v\) 是域 \(\mathrm{F}(t)\) 的元素，且

\[
\sum_ {j = 1} ^ {n} c _ {j} \frac {u _ {j} ^ {\prime}}{u _ {j}} + v ^ {\prime}
\]

属于环 F{[}t{]} ，则必有 \(v \in F[t]\) 。此外，若 \(t \in F\) ，则必有 \(u_{j} \in F\) （对 \(1 \leqslant j \leqslant n\) ）；若 \(t' / t \in F\) ，则对每个 j 存在整数 \(v_{j} \geqslant 0\) 使得 \(u_{j} / t^{v_{j}} \in F\) 。（在 F 的一个适当的伽罗瓦扩张中，把 \(u_{j}' / u_{j}\) 与 v 分解为简单元素，并利用习题 26）。*

* 28) a) 设 \(\mathrm{F} \subset \mathrm{M}(\Delta)\) 是一个微分域。称域 \(\mathrm{F}' \supset \mathrm{F}\) 为 \(\mathrm{F}\) 的初等扩张，如果存在有限序列

\[
\mathrm{F} = \mathrm{F} _ {0} \subset \mathrm{F} _ {1} \subset \mathrm{F} _ {2} \subset \dots \subset \mathrm{F} _ {n - 1} \subset \mathrm{F} _ {n} = \mathrm{F} ^ {\prime}\tag{1}
\]

使得对每个 \(j \leqslant n - 1\) 有 \(\mathrm{F}_{j + 1} = \mathrm{F}_j(t_j)\) ，其中或者有 \(t_j' = a_j' / a_j\) （对某个 \(a_j \neq 0\) ，它属于 \(\mathrm{F}_j\) ；从而可写 \(t_j = \log a_j\) ），或者有 \(t_j' / t_j = a_j'\) （对某个 \(a_j \in \mathrm{F}_j\) ；从而 \(t_j = \exp(a_j)\) ）。初等函数是属于 \(\mathbf{C}(z)\) 的某个初等扩张的函数（即 \(\mathbf{C}\) 上的有理函数域）。

例如，函数

\[
\operatorname{Arc} \sin z, \quad (\log z) ^ {\log z}, \quad \left(1 - \exp \left(\frac {1}{\exp \left(\frac {1}{z}\right) - 1}\right)\right) ^ {\alpha} \quad (\alpha \in \mathbf {C})
\]

都是初等函数。

\begin{enumerate}
\def\labelenumi{\alph{enumi})}
\setcounter{enumi}{1}
\tightlist
\item
  设 \(a \in \mathbf{F}\) ，使得存在元素 \(u_{j}\) （ \(1 \leqslant j \leqslant m\) ）与 \(v\) （ \(\mathrm{F}'\) 是 \(\mathrm{F}\) 的一个初等扩张），以及常数 \(\gamma_{j} \in \mathbf{C}\) （ \(1 \leqslant j \leqslant m\) ），满足
\end{enumerate}

\[
a = \sum_ {j = 1} ^ {m} \gamma_ {j} \frac {u _ {j} ^ {\prime}}{u _ {j}} + v ^ {\prime}.\tag{*}
\]

试证：这时存在元素 \(f_{j}\) （ \(1 \leqslant j \leqslant p\) ）与 \(g\) （在 \(F\) 中），以及常数 \(\beta_{j} \in \mathbf{C}\) （ \(1 \leqslant j \leqslant p\) ），使得

\[
a = \sum_ {j = 1} ^ {p} \beta_ {j} \frac {f _ {j} ^ {\prime}}{f _ {j}} + g ^ {\prime}.\tag{**}
\]

（对序列 (1) 的长度 \(n\) 作归纳，化归到 \(\mathrm{F}' = \mathrm{F}(t)\) 的情形。先修改 \(u_{j}\) ，证明可设诸 \(\gamma_{j}\) 在 \(\mathbf{Q}\) 上线性无关。当 \(t\) 在 \(\mathrm{F}\) 上超越时，利用习题 27：若 \(t' = s'/s\) （ \(s \in \mathrm{F}\) ），则必有 \(u_{j} \in \mathrm{F}\) 且 \(c \in \mathrm{F}[t]\) ；利用 (*) 并用反证法证明必有 \(v = \alpha t + b\) （ \(\alpha \in \mathbf{C}\) ， \(b \in \mathrm{F}\) ）。若 \(t'/t = r'\) （ \(r \in \mathrm{F}\) ），试证：把 \(v\) 换成 \(v + k.r\) （对适当的整数 \(k \in \mathbf{Z}\) ）后，可设 \(u_{j} \in \mathrm{F}\) ， \(v \in \mathrm{F}[t]\) ；用反证法证明 \(v\) 的次数为 0 ，故 \(v \in \mathrm{F}\) 。最后，若 \(t\) 在 \(\mathrm{F}\) 上代数，则把 \(\mathrm{F}'\) 嵌入某个伽罗瓦扩张 \(\mathrm{K}\) （ \(\mathrm{F}\) 的），并考虑 (*) 经 \(\mathrm{K}\) 在 \(\mathrm{F}.)\) 上的伽罗瓦群作用得到的变换。）*

* 29) 设 \(f\) 与 \(g\) 是 \(z\) 的两个有理函数（ \(\mathbf{C}(z)\) 中的元素）。试证：为使原函数 \(\int f(z)e^{g(z)}dz\) 是初等函数（习题 28），必须且只需存在有理函数 \(r\in \mathbf{C}(z)\) 使得 \(f = r' + rg'\) 。（令 \(t = e^g\) ，并考虑微分域 \(F = C(z,t)\) ；它是 \(\mathbf{C}(z)\) 的一个初等扩张，因为 \(t' / t = g'\) 。首先证明 \(t\) 在 \(\mathbf{C}(z)\) 上超越；用反证法，考虑一个伽罗瓦扩张 \(K\) （在 \(\mathbf{C}(z)\) 之上且包含 \(t\) ），以及方程 \(t' / t = g'\) 在 \(K\) 的伽罗瓦群作用下的变换；将得到方程 \(g' = u' / u\) （ \(u\in \mathbf{C}(z)\) ），而 \(g'\) 在 \(\mathbf{C}\) 中不可能有非单极点的极点。应用习题 28 ，有 \(ft = \sum_{j}\gamma_{j}\frac{u_{j}'}{u_{j}} +v'\) ，其中 \(\gamma_{j}\in \mathbf{C}\) ，而 \(u_{j}\) 与 \(v\) 属于 \(F\) ；可以化归到下述情形：

其中诸 \(\gamma_{j}\) 在 Q 上线性无关。然后，对扩张 \(\mathbf{C}(z)(t)\) 应用习题 27 ，证明有 \(ft = v' + h\) ，其中 \(h \in \mathbf{C}(z)\) ，\(v \in \mathbf{C}(z)[t]\) 。由此得出：v 必为 t 的一次多项式，从而 f 等于 \(v'\) 中 t 的系数。

由此推出：原函数 \(\int e^{z^2}dz\) 与 \(\int e^z dz / z\) 都不是初等函数（刘维尔定理），办法是在关系式 \(f = r' - rg'\) 中考察两边在 \(r\) 的一个极点的邻域内的性态。*

\begin{enumerate}
\def\labelenumi{\arabic{enumi})}
\setcounter{enumi}{29}
\tightlist
\item
  \begin{enumerate}
  \def\labelenumii{\alph{enumii})}
  \tightlist
  \item
    若 \(m\) 与 \(n\) 是满足 \(0 < m < n\) 的两个整数，试证
  \end{enumerate}
\end{enumerate}

\[
\int_ {0} ^ {+ \infty} \frac {x ^ {m - 1}}{1 + x ^ {n}} d x = \frac {\pi}{n \sin \frac {m \pi}{n}}.
\]

\begin{enumerate}
\def\labelenumi{\alph{enumi})}
\setcounter{enumi}{1}
\tightlist
\item
  试证：当 \(0 < a < 1\) 时，积分 \(\int_0^{+\infty} \frac{x^{a-1}}{1 + x} dx\) 对在紧区间中变动的 \(a\) 一致收敛，并由 \(a\) ) 推出
\end{enumerate}

\[
\int_ {0} ^ {+ \infty} \frac {x ^ {a - 1}}{1 + x} d x = \frac {\pi}{\sin a \pi}.
\]

\begin{enumerate}
\def\labelenumi{\arabic{enumi})}
\setcounter{enumi}{30}
\tightlist
\item
  若 \(I_{m,n}\) 是 \(\sin^{m}x\cos^{n}x\) 的一个原函数（m 与 n 为任意实数），试证
\end{enumerate}

\[
\mathrm{I} _ {m + 2, n} = - \frac {\sin^ {m + 1} x \cos^ {n + 1} x}{m + n + 2} + \frac {m + 1}{m + n + 2} \mathrm{I} _ {m, n}
\]

是 \(\sin^{m+2}x\cos^{n}x\) 的一个原函数（当 \(m+n+2\neq0\) 时）。

借助这个公式重新得到 III, p.~100 的公式 (29) ，并证明公式

\[
\int_ {0} ^ {\pi / 2} \sin^ {2 n + 1} x d x = \frac {2 . 4 . 6 \dots 2 n}{1 . 3 . 5 \dots (2 n + 1)} \quad (n \text {   an   integer   } \geqslant 0).
\]

\begin{enumerate}
\def\labelenumi{\arabic{enumi})}
\setcounter{enumi}{31}
\tightlist
\item
  证明沃利斯公式
\end{enumerate}

\[
\lim _ {n \rightarrow \infty} \frac {1}{\sqrt {n}}. \frac {2 . 4 . 6 \dots 2 n}{1 . 3 . 5 \dots (2 n - 1)} = \sqrt {\pi}
\]

试利用习题 31 及不等式 \(\sin^{n+1}x \leqslant \sin^{n}x\) （对 \(0 \leqslant x \leqslant \pi/2\) ）。

\begin{enumerate}
\def\labelenumi{\arabic{enumi})}
\setcounter{enumi}{32}
\tightlist
\item
  \begin{enumerate}
  \def\labelenumii{\alph{enumii})}
  \tightlist
  \item
    计算积分
  \end{enumerate}
\end{enumerate}

\[
\int_ {0} ^ {1} \left(1 - x ^ {2}\right) ^ {n} d x, \quad \int_ {0} ^ {+ \infty} \frac {d x}{\left(1 + x ^ {2}\right) ^ {n}}
\]

其中 n 是整数 \textgreater{} 0 ，借助习题 31 。

\begin{enumerate}
\def\labelenumi{\alph{enumi})}
\setcounter{enumi}{1}
\tightlist
\item
  试证：
\end{enumerate}

\[
\begin{array}{l l} 1 - x ^ {2} \leqslant e ^ {- x ^ {2}} & \text { for } 0 \leqslant x \leqslant 1 \\ e ^ {- x ^ {2}} \leqslant \frac {1}{1 + x ^ {2}} & \text { for } x \geqslant 0. \end{array}
\]

\begin{enumerate}
\def\labelenumi{\alph{enumi})}
\setcounter{enumi}{2}
\tightlist
\item
  由 \(a\) 与 \(b\) ) 以及沃利斯公式（习题 32）推出
\end{enumerate}

\[
\int_ {0} ^ {+ \infty} e ^ {- x ^ {2}} d x = \frac {\sqrt {\pi}}{2}.
\]

\begin{enumerate}
\def\labelenumi{\arabic{enumi})}
\setcounter{enumi}{33}
\tightlist
\item
  \begin{enumerate}
  \def\labelenumii{\alph{enumii})}
  \tightlist
  \item
    试证：对 \(\alpha > 0\) ，下列函数的导数
  \end{enumerate}
\end{enumerate}

\[
\mathrm{I} (\alpha) = \int_ {0} ^ {+ \infty} e ^ {- \alpha x} \frac {\sin x}{x} d x
\]

等于 \(-\int_{0}^{+\infty}e^{-\alpha x}\sin x dx\) 。

\begin{enumerate}
\def\labelenumi{\alph{enumi})}
\setcounter{enumi}{1}
\tightlist
\item
  由此推出 \(\int_0^{+\infty}\frac{\sin x}{x} dx = \frac{\pi}{2}.\)
\end{enumerate}

\begin{enumerate}
\def\labelenumi{\arabic{enumi})}
\setcounter{enumi}{34}
\tightlist
\item
  通过对参数求导并利用习题 33 c) ，证明公式
\end{enumerate}

\[
\begin{array}{l} \int_ {0} ^ {+ \infty} e ^ {- x ^ {2}} \cos \alpha x d x = \frac {\sqrt {\pi}}{2} e ^ {- \alpha^ {2}}, \quad \int_ {0} ^ {+ \infty} \frac {1 - e ^ {- \alpha x ^ {2}}}{x ^ {2}} d x = \sqrt {\pi \alpha} \\ \int_ {0} ^ {+ \infty} \exp \left(- x ^ {2} - \frac {\alpha^ {2}}{x ^ {2}}\right) d x = \frac {\sqrt {\pi}}{2} e ^ {- 2 \alpha} \quad (\alpha > 0). \end{array}
\]

\begin{enumerate}
\def\labelenumi{\arabic{enumi})}
\setcounter{enumi}{35}
\tightlist
\item
  由上面的习题 33 c) 以及 II, p.~88, exerc. 9 推出
\end{enumerate}

\[
\int_ {0} ^ {+ \infty} \sin x ^ {2} d x = \frac {1}{2} \sqrt {\frac {\pi}{2}}.
\]

\begin{enumerate}
\def\labelenumi{\arabic{enumi})}
\setcounter{enumi}{36}
\tightlist
\item
  设 \(\mathbf{f}\) 是 \(]0, 1[\) 上的一个规则向量函数，使得积分 \(\int_0^\pi \mathbf{f}(\sin x) dx\) 收敛。试证：积分 \(\int_0^\pi x \mathbf{f}(\sin x) dx\) 收敛，并且
\end{enumerate}

\[
\int_ {0} ^ {\pi} x \mathbf {f} (\sin x) d x = \frac {\pi}{2} \int_ {0} ^ {\pi} \mathbf {f} (\sin x) d x.
\]

\begin{enumerate}
\def\labelenumi{\arabic{enumi})}
\setcounter{enumi}{37}
\item
  设 \(\mathbf{f}\) 是对 \(x \geqslant 0\) 定义的规则向量函数，在点 \(x = 0\) 处连续，且使得积分 \(\int_{a}^{+\infty} \mathbf{f}(x) dx / x\) 对 \(a > 0\) 收敛。试证：对 \(a > 0\) 与 \(b > 0\) ，积分 \(\int_{0}^{+\infty} \frac{\mathbf{f}(ax) - \mathbf{f}(bx)}{x} dx\) 收敛且等于 \(\mathbf{f}(0) \log a / b\) 。
\item
  设 \(m\) 是 \([0, +\infty[\) 上的一个凸函数，满足 \(m(0) = 0\) ，并设 \(p\) 是满足 \(-1 < p < +\infty\) 的数。试证：积分 \(\int_0^{+\infty} x^p \exp(-m_r'(x)) dx\) 收敛，积分 \(\int_0^{+\infty} x^p \exp(-m(x)/x) dx\) 也收敛，并且
\end{enumerate}

\[
\int_ {0} ^ {+ \infty} x ^ {p} \exp (- m (x) / x) d x \leqslant e ^ {p + 1} \int_ {0} ^ {+ \infty} x ^ {p} \exp (- m _ {r} ^ {\prime} (x)) d x.
\]

（对 \(k > 1\) 与 \(A > 0\) ，注意 \(m(kx) \geqslant m(x) + (k - 1)x m_r'(x)\) ，并由此推出不等式

\[
k ^ {- p - 1} \int_ {0} ^ {k A} x ^ {p} \exp (- m (x) / x) d x \leqslant \int_ {0} ^ {A} x ^ {p} \exp \left(- \frac {m (x)}{k x} - \frac {k - 1}{k} m _ {r} ^ {\prime} (x)\right) d x.
\]

利用赫尔德不等式（III, p.~115, exerc. 3）估计第二个积分，然后让 A 趋于 \(+\infty\) 并让 \(k\) 趋于 1。）

\setcounter{exercisegroup}{1}
\refstepcounter{exercisegroup}
\section*{第 \S\arabic{exercisegroup} 节习题}
\phantomsection
\addcontentsline{toc}{section}{第 \protect\S\arabic{exercisegroup} 节习题}

\begin{enumerate}
\def\labelenumi{\arabic{enumi})}
\tightlist
\item
  设 \(\mathbf{f}\) 是一个 \(n\) 次可导的向量函数（在区间 \(\mathrm{I} \subset \mathbf{R}\) 上）。证明公式
\end{enumerate}

\[
\mathrm{D} ^ {n} \mathbf {f} \left(e ^ {x}\right) = \sum_ {m = 1} ^ {n} \frac {a _ {m}}{m !} e ^ {m x} \mathbf {f} ^ {(m)} \left(e ^ {x}\right)
\]

在使 \(e^{x} \in I\) 的每一点 x 处成立，其中系数 \(a_{m}\) 可表为

\[
a _ {m} = \sum_ {p = 0} ^ {m} (- 1) ^ {p} \binom {m} {p} (m - p) ^ {n}
\]

（用 I, p.~41, exerc. 7 的方法，并利用 \(e^{x}\) 的泰勒展开）。

\begin{enumerate}
\def\labelenumi{\arabic{enumi})}
\setcounter{enumi}{1}
\item
  设 \(f\) 是一个实函数，在一点 \(n\) 次可导，该点为 \(x\) ；\(\mathbf{g}\) 是一个向量函数，\(n\) 次可导，可导点为 \(f(x)\) 。若令 \(\mathrm{D}^n(\mathbf{g}(f(x))) = \sum_{k=1}^{n} \mathbf{g}^{(k)}(f(x)) u_k(x)\) ，则 \(u_k\) 仅依赖于函数 \(f\) ；由此推出：\(u_k(x)\) 是 \(t^k\) 的系数（在 \(e^{-tf(x)}\mathrm{D}^n(e^{tf(x)})\) 按 \(t\) 展开的展开式中）。
\item
  对每个实数 \(x > 0\) 与每个复数 \(m = \mu + iv\) ，令 \(x^{m} = e^{m\log x}\) ；试证：III, p.~108 的公式 (19) 对复的 \(m\) 与 \(x > -1\) 仍然有效，且该公式中的余项 \(r_{n}(x)\) 满足不等式
\end{enumerate}

\[
\begin{array}{l l} | r _ {n} (x) | \leqslant \left| \binom {m - 1} {n} \frac {m}{\mu} x ^ {n} \Big [ (1 + x) ^ {\mu} - 1 \Big ] \right| & \text { if } \mu \neq 0 \\ | r _ {n} (x) | \leqslant \left| \binom {m - 1} {n} m x ^ {n} \log (1 + x) \right| & \text { if } \mu = 0. \end{array}
\]

把二项级数收敛性的研究推广到 m 为复数的情形。

\begin{enumerate}
\def\labelenumi{\arabic{enumi})}
\setcounter{enumi}{3}
\tightlist
\item
  对每个实数 \(x\) 及每个数 \(p > 1\) ，证明不等式
\end{enumerate}

\[
\begin{array}{r l} | 1 + x | ^ {p} \leqslant 1 + p x + \frac {p (p - 1)}{2} x ^ {2} + \dots + \frac {p (p - 1) \dots (p - m + 2)}{(m - 1) !} x ^ {m - 1} \\ & + \frac {p (p - 1) \dots (p - m + 1)}{m !} | x | ^ {m} + h _ {p} | x | ^ {p} \end{array}
\]

其中令 \(m = [p]\) （即 \(p\) 的整数部分），并且

\[
h _ {p} = \frac {p (p - 1) \dots (p - m + 1)}{(m - 1) !} \int_ {0} ^ {1} z ^ {p - m} (1 - z) ^ {m - 1} d z.
\]

¶ 5) 试证：\(\pi\) 是无理数，方法如下：若 \(\pi = p/q\) （ \(p\) 与 \(q\) 为整数），则令 \(f(x) = (x(\pi - x))^n / n!\) 后，积分 \(q^n \int_0^\pi f(x) \sin x dx\) 将是一个整数 \(>0\) （利用 \(n + 1\) 阶分部积分公式）；但试证另一方面，\(q^n \int_0^\pi f(x) \sin x dx\) 当 \(n\) 趋于 \(+\infty\) 时趋于 0 。

\begin{enumerate}
\def\labelenumi{\arabic{enumi})}
\setcounter{enumi}{5}
\tightlist
\item
  试证：在区间 \([-1,+1]\) 上，函数 \(|x|\) 是多项式的一致极限，方法是注意 \(|x|=(1-(1-x^{2}))^{1/2}\) 并利用二项级数。由此得出魏尔斯特拉斯定理的另一证明（cf.~II, p.~83, exerc. 20）。
\end{enumerate}

¶ 7) 设 \(p\) 是素数，\(\mathbf{Q}_p\) 是 \(p\) 进数域（Gen.~Top., III, p.~322 and 323, exercs. 23 to 25），\(\mathbf{Z}_p\) 是 \(p\) 进整数环，\(\mathfrak{p}\) 是主理想 \((p)\) （在 \(\mathbf{Z}_p\) 中）。

\begin{enumerate}
\def\labelenumi{\alph{enumi})}
\tightlist
\item
  设 \(a = 1 + pb\) ，其中 \(b \in \mathbf{Z}_p\) 是乘法群 \(1 + \mathfrak{p}\) 的元素；试证：当有理整数 \(m\) 无限增大时， \(p\) 进数 \(\frac{(1 + pb)^{p^m} - 1}{p^m}\) 趋于一个极限，它等于收敛级数
\end{enumerate}

\[
\frac {p b}{1} - \frac {p ^ {2} b ^ {2}}{2} + \dots + (- 1) ^ {n - 1} \frac {p ^ {n} b ^ {n}}{n} + \dots
\]

把这个极限记作 \(\log a\) 。

\begin{enumerate}
\def\labelenumi{\alph{enumi})}
\setcounter{enumi}{1}
\item
  试证：当 \(p\) 进数 \(x\) 在 \(\mathbf{Q}_p\) 中趋于 0 时，数 \(\frac{a^x - 1}{x}\) 趋于 \(\log a\) （利用 \(a\) ) 及 \(\mathbf{Q}_p\) 的拓扑的定义）。
\item
  试证：若 \(p \neq 2\) ，则 \(\log a \equiv pb \pmod{\mathfrak{p}^2}\) ；若 \(p = 2\) ，则 \(\log a \equiv 0 \pmod{\mathfrak{p}^2}\) ，且 \(\log a \equiv -4b^4 \pmod{\mathfrak{p}^3}\) 。
\item
  试证：若 \(p \neq 2\) （相应地，p = 2），则 \(x \mapsto \log x\) 是乘法拓扑群 \(1 + p\) 到加法拓扑群 p 上的同构（相应地，到 \(p^{2}\) 上）；特别地，若 \(e_{p}\) 是 \(1 + p\) 中使 \(\log e_{p} = p\) 的元素（相应地，\(\log e_{2} = 4\) ），则 \(Z_{p}\) 到 \(1 + p\) 上的、与 \(x \mapsto \frac{1}{p} \log x\) 互逆的同构（相应地，与 \(x \mapsto \frac{1}{4} \log x\) 互逆者）就是 \(y \mapsto e_{p}^{y}\) （cf.~Gen.~Top., III, p.~323, exerc. 25）。
\item
  试证：对每个 \(a \in 1 + p\) ，连续函数 \(x \mapsto a^{x}\) （定义在 \(Z_{p}\) 上）在每一点处的导数等于 \(a^{x} \log a\) ；由此推出函数 \(\log x\) 在 \(1 + p\) 的每一点处的导数等于 1/x 。
\end{enumerate}

¶ 8) a) 用习题 7 的记号，试证：通项为 \(x^n / n!\) 的级数对每个 \(x \in \mathfrak{p}\) 收敛（若 \(p \neq 2\) ），对每个 \(x \in \mathfrak{p}^2\) （但对任何 \(x \notin \mathfrak{p}^2\) 都不）收敛（若 \(p = 2\) ）（定出 \(p\) 在 \(n!\) 的素因子分解中的指数）。若 \(f(x)\) 是这级数的和，试证 \(f\) 是 \(\mathfrak{p}\) （相应地，\(\mathfrak{p}^2\) ）到 \(1 + \mathfrak{p}^2\) 中的连续同态。由此推出：对每个 \(z \in \mathbf{Z}_p\) ，有 \(f(pz) = e_p^z\) （相应地，\(f(p^2z) = e_p^z\) ）；换言之，

\[
e _ {p} = 1 + \frac {p}{1 !} + \frac {p ^ {2}}{2 !} + \dots + \frac {p ^ {n}}{n !} + \dots \quad \text { if } p \neq 2
\]

\[
e _ {2} = 1 + \frac {4}{1 !} + \frac {4 ^ {2}}{2 !} + \dots + \frac {4 ^ {n}}{n !} + \dots
\]

（利用习题 7 e)）。

\begin{enumerate}
\def\labelenumi{\alph{enumi})}
\setcounter{enumi}{1}
\tightlist
\item
  对每个 \(a \in 1 + \mathfrak{p}\) 及每个 \(x \in \mathbf{Z}_p\) ，试证 \(\log (a^x) = x\log a\) ，并由 a) 与习题 7 d) 推出
\end{enumerate}

\[
a ^ {x} = 1 + \frac {x \log a}{1 !} + \frac {x ^ {2} (\log a) ^ {2}}{2 !} + \dots + \frac {x ^ {n} (\log a) ^ {n}}{n !} + \dots
\]

\begin{enumerate}
\def\labelenumi{\alph{enumi})}
\setcounter{enumi}{2}
\tightlist
\item
  对每个 \(m \in Z_{p}\) ，试证：连续函数 \(x \mapsto x^{m}\) （定义在 \(1 + p\) 上）的导数等于 \(mx^{m-1}\) （利用 b) 与习题 7 e)）。
\end{enumerate}

\(d\) ) 试证：对每个 \(m\in \mathbf{Z}_p\) 及每个 \(x\in \mathfrak{p}\) （当 \(p\neq 2\) ）或每个 \(x\in \mathfrak{p}^2\) （当 \(p = 2\) ），通项为 \(\binom{m}{n}x^n\) 的级数收敛，且其和是 \(m\) 的连续函数；由此推出这个和等于 \((1 + x)^m\) ，方法是注意 \(\mathbf{Z}\) 在 \(\mathbf{Z}_p\) 中（处处）稠密。

¶ 9) 用习题 7) 的记号，把 \(\mathbf{O}_2^+ (\mathbf{Q}_p)\) （空间 \(\mathbf{Q}_p^2\) 的旋转群）定义为形如

\[
\left( \begin{array}{c c} x & y \\ - y & x \end{array} \right)
\]

的矩阵所成的群，其元素属于 \(\mathbf{Q}_p\) ，满足 \(x^2 + y^2 = 1\) ，该群赋有 Gen.~Top., VIII, p.~125, exerc. 2 中定义的拓扑。

\begin{enumerate}
\def\labelenumi{\alph{enumi})}
\item
  用 \(G_{n}\) 表示 \(\mathbf{O}_{2}^{+}(\mathbf{Q}_{p})\) 中由满足 \(t=y/(1+x)\in\mathfrak{p}^{n}\) 的矩阵所成的子群。试证：\(G_{n}\) 是紧群，\(G_{n}/G_{n+1}\) 同构于 Z/pZ，且 \(G_{1}\) 的仅有的紧子群就是诸群 \(G_{n}\) （cf.~Gen.~Top., III, p.~323, exerc. 24）。
\item
  试证：\(G_{1}\) 与由矩阵
\end{enumerate}

\[
\left( \begin{array}{c c} x & y \\ - y & x \end{array} \right)
\]

所成的子群相同，其中 \(x^{2} + y^{2} = 1\) ，且 \(x \in 1 + p^{2}\) ， \(y \in p\) （若 \(p \neq 2\) ）；\(x \in 1 + p^{3}\) ， \(y \in p^{2}\) （若 p = 2）。

\begin{enumerate}
\def\labelenumi{\alph{enumi})}
\setcounter{enumi}{2}
\tightlist
\item
  试证：通项为 \((-1)^{n} x^{2n}/(2n)!\) 与 \((-1)^{n-1} x^{2n+1}/(2n+1)!\) 的级数对每个 \(x \in p\) 收敛（若 \(\neq 2\) ），对每个 \(x \in p^{2}\) 收敛（若 p = 2）；设 \(\cos x\) 与 \(\sin x\) 是这两个级数的和。试证映射
\end{enumerate}

\[
x \mapsto \left( \begin{array}{c c} \cos x & \sin x \\ - \sin x & \cos x \end{array} \right)
\]

是加法拓扑群 \(\mathfrak{p}\) （相应地，\(\mathfrak{p}^2\) ）到群 \(G_1\) 上的同构。

\(d\) ) 若 \(p\) 形如 \(4h + 1\) （ \(h\) 为整数），则 \(\mathbf{Q}_p\) 中存在元素 \(i\) 使得 \(i^2 = -1\) 。对每个 \(z \in \mathbf{Q}_p^*\) 关联矩阵

\[
\left( \begin{array}{c c} \frac {1}{2} \left(z + \frac {1}{z}\right) & \frac {1}{2 i} \left(z - \frac {1}{z}\right) \\ - \frac {1}{2 i} \left(z - \frac {1}{z}\right) & \frac {1}{2} \left(z + \frac {1}{z}\right) \end{array} \right)
\]

便定义了乘法群 \(Q_{p}^{*}\) 到群 \(\mathbf{O}_{2}^{+}(\mathbf{Q}_{p})\) 上的一个同构；在这个同构下，群 \(1 + p\) 对应于 \(G_{1}\) ，并且有 \(\cos px + i \sin px = e_{p}^{ix}\) （III, p.~126, exerc. 8）。

\begin{enumerate}
\def\labelenumi{\alph{enumi})}
\setcounter{enumi}{4}
\tightlist
\item
  若 p 形如 \(4h + 3\) （h 为整数），则 \(\mathbf{O}_{2}^{+}(\mathbf{Q}_{p})\) 的矩阵的元素必为 p 进整数。此时多项式 \(X^{2} + 1\) 在 \(Q_{p}\) 中不可约；设 \(\mathbf{Q}_{p}(i)\) 是 \(Q_{p}\) 的一个二次扩张，由添加 \(X^{2} + 1\) 的一个根 i 得到；给 \(\mathbf{Q}_{p}(i)\) 赋以 Gen.~Top., VIII, p.~127, exerc. 2 中定义的拓扑。群 \(\mathbf{O}_{2}^{+}(\mathbf{Q}_{p})\) 同构于 \(\mathbf{Q}_{p}(i)\) 中范数为 1 的元素所成的乘法群 N，同构的方式是把矩阵 \(\begin{pmatrix} x & y \\ -y & x \end{pmatrix}\) 对应于元素 \(z = x + iy\) 。试证：有 \(p + 1\) 个数满足方程 \(x^{p+1} = 1\) （在 \(\mathbf{Q}_{p}(i)\) 中），且它们构成 N 的一个循环子群 R（论证方法同 Gen.~Top., III, p.~323, exerc. 24；然后证明有 \(p + 1\) 个不同的根满足同余式 \(x^{p+1} \equiv 1 \pmod{p}\) （在 \(\mathbf{Q}_{p}(i)\) 中），并对该同余式的每个根 a 作序列 \((a^{p2n})\) ）。由此推出：群 \(\mathbf{O}_{2}^{+}(\mathbf{Q}_{p})\) 同构于群 R 与 \(G_{1}\) 的乘积。
\end{enumerate}

\(f\) ) 试证：当 \(p = 2\) 时，群 \(\mathbf{O}_2^+ (\mathbf{Q}_p)\) 同构于群 \(\mathrm{G}_1\) 与 4 阶循环群的乘积。
